\chapter{Literature Review}

\section{Impact Stress Distribution and Mechanics}
Impact stress analysis on construction helmet shells is a critical engineering 
domain. Under EN 397 conditions, shells experience peak stresses concentrated 
at the dome apex. According to Long et al. \cite{long2015}, finite element 
models applying explicit solvers are effective at mapping these transient 
loads. However, a primary limitation in such models is the frequent 
adoption of simplified linear elastic material assumptions for the underlying 
foam liners, which may fail to capture the energy dissipation associated with 
non-linear cell crushing. Research indicates that standard HDPE shells 
experience peak stresses of 2.55--3.57 kg/mm$^2$ with moderate deformation, while ABS 
shells undergo peak stresses of 3.06--4.59 kg/mm$^2$ due to their higher stiffness.

While Wu et al. \cite{wu2022} report that helmet shells remain primarily 
within elastic deformation limits (strains of 4\%--6.3\%), there is 
substantial uncertainty regarding material response under high-strain-rate 
impact loading. Specifically, the energy absorption capacity is fundamentally 
dictated by the material's elastic modulus; however, the transition from 
structural yielding in ductile materials to load distribution in stiff 
materials is often treated as a binary mechanism. This study identifies a 
critical gap in understanding how this transition influences transmitted 
impulse when sustainable, non-conventional polymers are introduced.

\section{Conventional Material Properties}
Conventional materials have dictated helmet design for decades, yet their 
comparative performance under varying impact velocities remains under-scrutinised.
\begin{itemize}
    \item \textbf{HDPE:} Widely used due to its semi-rigid nature and high 
    ductility. While industrial data suggests a yield strength of 
    approx.\ 2.55 kg/mm$^2$ \cite{cruz2023}, academic studies often 
    report conflicting results depending on the molecular weight and 
    processing conditions, indicating that standard HDPE performance may 
    vary measurably in field applications.
    \item \textbf{ABS:} Known for exceptional rigidity and tensile strength 
    (4.08--5.10 kg/mm$^2$). While Author A reports high energy absorption 
    in HDPE, Author B suggests reduced performance under transient loading, 
    indicating uncertainty in how flexural modulus (approx. 214.14 kg/mm$^2$) 
    interacts with the localized buckling resistance required for 
    EN 397 compliance.
\end{itemize}

\section{Sustainable and Environmentally Friendly Materials}
Environmental sustainability considerations are driving helmet material 
innovation, but existing research is primarily restricted to 
quasi-static contexts, which fail to represent the high-strain-rate 
environment of an industrial impact.
\begin{itemize}
    \item \textbf{Bio-Based Polymers:} Materials such as Bio-PE offer 
    potential impact resistance equivalent to fossil-derived PE \cite{alvarez2022}. 
    However, the lack of documented high-velocity impact data for these 
    renewable biomass feedstocks represents a major hurdle for regulatory 
    acceptance in safety-critical PPE.
    \item \textbf{Natural Fibre Composites:} Integrating bamboo or coir 
    fibres into thermoplastic matrices increases tensile strength 
    \cite{obele2015, chandramohan2015}. Research by Chen et al. \cite{chen2020} 
    suggests 10\% fibre loading is optimal; however, their study 
    focuses on tensile rather than impact dynamics. The structural 
    viability of higher loadings (up to 30\%) remains contested, 
    particularly regarding the interfacial bonding required to maintain 
    integrity under transient EN 397 loads.
    \item \textbf{Recycled Polycarbonate (rPC):} Studies indicate that rPC 
    can maintain performance parity with virgin filaments \cite{baechler2013}. 
    However, these studies often overlook the effects of polymer 
    degradation during multiple recycling cycles, which can induce 
    brittleness. While its extreme rigidity promotes reduced tangential 
    frictional interaction \cite{mills2011}, the potential for catastrophic 
    brittle fracture under peak EN 397 loads remains an unaddressed risk 
    in current literature.
\end{itemize}

While explicit FEA has benchmarked conventional safety helmets \cite{long2015, wu2022}, 
sustainable alternatives lack a validated high-velocity performance baseline. 
This investigation addresses this critical deficiency by establishing a 
comparative framework to quantify their kinetics and regulatory compliance 
under EN 397 transient loading, explicitly accounting for the uncertainties 
identified in existing literature.
